% ============================================================
\chapter{Reproducibility in Research --- Standards and Best Practices}
\label{ch:reproducibility}
\index{reproducibility}
% ============================================================

\section{The Reproducibility Crisis in ML}
\label{sec:repro-crisis}

A growing body of evidence shows that many published ML results cannot
be reproduced. Surveys report that over 60\% of ML experiments lack
sufficient detail for independent verification. The causes are
multiple: missing code, undocumented hyperparameters, unreported
hardware configurations, and non-deterministic training procedures.

\begin{definition}[Reproducibility]
\textbf{Reproducibility}\index{reproducibility!definition} is the
ability for an independent team to obtain the same results using the
\emph{same} code and data. It is distinct from \emph{replicability},
which aims to confirm findings using \emph{different} code or data.
\end{definition}

\section{Levels of Reproducibility}
\label{sec:repro-levels}
\index{reproducibility!levels}

The ACM distinguishes three levels with increasing stringency:

\begin{center}
\begin{tabular}{@{}llp{6.5cm}@{}}
\toprule
\textbf{Level} & \textbf{Definition} & \textbf{Requirements} \\
\midrule
Repeatability & Same team, same setup & Fixed seeds, versioned code,
  deterministic pipeline \\
Reproducibility & Different team, same artifacts & Published code,
  data, environment specs, clear instructions \\
Replicability & Different team, different setup & Detailed methodology,
  model cards, independent confirmation \\
\bottomrule
\end{tabular}
\end{center}

\begin{center}
\begin{tikzpicture}[
  level/.style={rectangle, draw, rounded corners, minimum width=10cm,
    minimum height=1.2cm, align=center, font=\small},
  arr/.style={-{Stealth[length=2.5mm]}, thick, steelblue}
]
  \node[level, fill=green!15] (rep) {\textbf{Repeatability}\\
    Same team obtains same results on same hardware};
  \node[level, fill=yellow!15, above=0.5cm of rep] (repro)
    {\textbf{Reproducibility}\\
    Different team obtains same results with same code/data};
  \node[level, fill=orange!15, above=0.5cm of repro] (repl)
    {\textbf{Replicability}\\
    Different team confirms findings with different implementation};

  \draw[arr] (rep) -- node[right, font=\scriptsize] {harder} (repro);
  \draw[arr] (repro) -- node[right, font=\scriptsize] {hardest} (repl);
\end{tikzpicture}
\end{center}

\section{Random Seed Management}
\label{sec:random-seeds}
\index{random seed}

Non-determinism is the primary obstacle to repeatability. ML pipelines
have multiple sources of randomness that must all be controlled.

\begin{codeblock}[title=\textbf{Complete random seed management}]
\begin{minted}{python}
"""Full seed management for reproducibility."""
import os
import random

import numpy as np
import torch


def set_seed(seed: int = 42) -> None:
    """Set all random seeds for reproducibility.

    Controls randomness in:
    - Python's random module
    - NumPy
    - PyTorch (CPU and CUDA)
    - CUDA convolution algorithms
    - Hash-based operations
    """
    # Python
    random.seed(seed)
    os.environ["PYTHONHASHSEED"] = str(seed)

    # NumPy
    np.random.seed(seed)

    # PyTorch CPU
    torch.manual_seed(seed)

    # PyTorch CUDA (all GPUs)
    torch.cuda.manual_seed(seed)
    torch.cuda.manual_seed_all(seed)

    # Deterministic algorithms (may reduce performance)
    torch.backends.cudnn.deterministic = True
    torch.backends.cudnn.benchmark = False

    # PyTorch 2.0+ deterministic mode
    torch.use_deterministic_algorithms(True)
    os.environ["CUBLAS_WORKSPACE_CONFIG"] = ":4096:8"


# Use at the very start of training
set_seed(42)
\end{minted}
\end{codeblock}

\begin{warning}[title=\textbf{Limits of determinism on GPU}]
Even with all seeds fixed, some GPU operations remain non-deterministic
by default:
\begin{itemize}
  \item Atomic operations in CUDA (e.g., \texttt{scatter\_add})
  \item cuDNN auto-tuner selects different algorithms per run
  \item Multi-threaded data loading with \pkg{DataLoader}
    (\texttt{num\_workers > 0})
\end{itemize}
Setting \texttt{torch.use\_deterministic\_algorithms(True)} forces
deterministic alternatives but may slow down training by 10--20\%.
\end{warning}

\begin{codeblock}[title=\textbf{Deterministic DataLoader}]
\begin{minted}{python}
"""Reproducible data loading."""
import torch
from torch.utils.data import DataLoader


def seed_worker(worker_id: int) -> None:
    """Seed each DataLoader worker for reproducibility."""
    worker_seed = torch.initial_seed() % 2**32
    np.random.seed(worker_seed)
    random.seed(worker_seed)


generator = torch.Generator()
generator.manual_seed(42)

train_loader = DataLoader(
    dataset,
    batch_size=32,
    shuffle=True,
    num_workers=4,
    worker_init_fn=seed_worker,
    generator=generator,
)
\end{minted}
\end{codeblock}

\section{Model Cards}
\label{sec:model-cards}
\index{model card}

A \textbf{model card}\index{model card} (Mitchell et al., 2019) is a
standardized document that accompanies a trained model, providing
transparency about its intended use, performance, and limitations.

\begin{keyformulas}[title=\textbf{Model card template}]
\begin{enumerate}
  \item \textbf{Model Details}: architecture, version, training date,
    framework, license.
  \item \textbf{Intended Use}: primary use cases, out-of-scope uses.
  \item \textbf{Training Data}: dataset description, size,
    preprocessing, known biases.
  \item \textbf{Evaluation Data}: test set description, evaluation
    protocol.
  \item \textbf{Metrics}: accuracy, F1, AUC---disaggregated by
    demographic group.
  \item \textbf{Ethical Considerations}: potential harms, biases,
    fairness analysis.
  \item \textbf{Caveats and Recommendations}: known limitations,
    failure modes.
\end{enumerate}
\end{keyformulas}

\begin{codeblock}[title=\textbf{Model card generation with Python}]
\begin{minted}{python}
"""Generate a model card as structured YAML."""
import yaml
from datetime import date

model_card = {
    "model_details": {
        "name": "Fraud Detection Classifier v2.1",
        "architecture": "XGBoost",
        "framework": "xgboost==2.0.3",
        "training_date": str(date.today()),
        "license": "MIT",
    },
    "intended_use": {
        "primary": "Detect fraudulent credit card transactions",
        "out_of_scope": [
            "Non-financial fraud detection",
            "Real-time streaming (latency > 100ms)",
        ],
    },
    "training_data": {
        "dataset": "Internal transactions 2022-2024",
        "size": "2.4M transactions",
        "fraud_rate": "1.2%",
        "preprocessing": "SMOTE oversampling, standard scaling",
    },
    "metrics": {
        "overall": {"accuracy": 0.987, "f1": 0.912, "auc": 0.975},
        "by_region": {
            "north_america": {"f1": 0.923},
            "europe": {"f1": 0.908},
            "asia": {"f1": 0.891},
        },
    },
    "limitations": [
        "Lower performance on transactions > $10,000",
        "Not validated for cryptocurrency transactions",
        "Requires retraining quarterly due to drift",
    ],
}

with open("model_card.yaml", "w") as f:
    yaml.dump(model_card, f, default_flow_style=False, sort_keys=False)
\end{minted}
\end{codeblock}

\section{Datasheets for Datasets}
\label{sec:datasheets}
\index{datasheet}

\textbf{Datasheets for Datasets}\index{datasheet} (Gebru et al., 2021)
standardize dataset documentation. Key sections:

\begin{toolbox}[title=\textbf{Datasheet sections}]
\begin{description}
  \item[Motivation] Why was the dataset created? Who funded it?
  \item[Composition] What do instances represent? How many? Any
    missing data?
  \item[Collection process] How was data acquired? Consent process?
  \item[Preprocessing] What cleaning or transformations were applied?
  \item[Uses] Intended tasks? Tasks to avoid?
  \item[Distribution] How is the dataset shared? License?
  \item[Maintenance] Who maintains it? Update frequency? Deprecation
    policy?
\end{description}
\end{toolbox}

\section{NeurIPS Reproducibility Checklist}
\label{sec:neurips-checklist}
\index{NeurIPS checklist}

Major ML conferences now require authors to complete a reproducibility
checklist. The NeurIPS checklist includes:

\begin{bestpractice}[title=\textbf{NeurIPS-style reproducibility checklist}]
\begin{enumerate}
  \item Code is submitted (or will be released upon acceptance).
  \item All hyperparameters are reported, including search ranges.
  \item Training and evaluation datasets are described or cited.
  \item Compute resources used are documented (GPU type, hours).
  \item Number of runs and variance/confidence intervals are reported.
  \item Random seeds are specified.
  \item Data preprocessing steps are fully described.
  \item Software dependencies are listed with versions.
  \item Statistical significance tests are included where relevant.
  \item Negative results and failure cases are discussed.
\end{enumerate}
\end{bestpractice}

\begin{remark}
Reporting the mean and standard deviation over multiple runs (e.g., 5
runs with different seeds) is far more informative than a single best
result. A model with $92.1 \pm 0.3\%$ accuracy is more trustworthy
than one that claims $93.0\%$ from a single lucky seed.
\end{remark}

\section{Practical Reproducibility Toolkit}
\label{sec:repro-toolkit}

\begin{codeblock}[title=\textbf{Reproducibility configuration}]
\begin{minted}{yaml}
# configs/experiment.yaml
experiment:
  name: "fraud-detection-v2"
  seed: 42
  deterministic: true

data:
  path: "data/transactions.csv"
  dvc_hash: "a1b2c3d4"  # DVC file hash for traceability
  test_size: 0.2
  stratify: "target"

model:
  type: "xgboost"
  params:
    n_estimators: 500
    max_depth: 8
    learning_rate: 0.05
    random_state: 42

training:
  n_runs: 5  # Multiple runs for confidence intervals
  seeds: [42, 123, 456, 789, 1024]

environment:
  python: "3.11.7"
  gpu: "NVIDIA A100"
  cuda: "12.1"
  requirements_hash: "sha256:e5f6..."
\end{minted}
\end{codeblock}

\begin{codeblock}[title=\textbf{Multi-run experiment with statistics}]
\begin{minted}{python}
"""Run experiment multiple times for robust evaluation."""
import json
import numpy as np
from scipy import stats


def run_experiment(seed: int, config: dict) -> dict:
    """Run a single experiment with a given seed."""
    set_seed(seed)
    # ... training code ...
    return {"accuracy": acc, "f1": f1, "auc": auc}


seeds = [42, 123, 456, 789, 1024]
results = [run_experiment(s, config) for s in seeds]

# Aggregate results
for metric in ["accuracy", "f1", "auc"]:
    values = [r[metric] for r in results]
    mean = np.mean(values)
    std = np.std(values, ddof=1)
    ci = stats.t.interval(
        0.95, len(values) - 1,
        loc=mean, scale=stats.sem(values),
    )
    print(f"{metric}: {mean:.4f} +/- {std:.4f} "
          f"(95% CI: [{ci[0]:.4f}, {ci[1]:.4f}])")
\end{minted}
\end{codeblock}

\section{Best Practices and Anti-patterns}

\begin{bestpractice}[title=\textbf{Reproducibility best practices}]
\begin{enumerate}
  \item \textbf{Fix all seeds}: Python, NumPy, PyTorch, CUDA, data
    loaders.
  \item \textbf{Version everything}: code (Git), data (DVC),
    experiments (MLflow), environments (Docker).
  \item \textbf{Report multiple runs}: mean $\pm$ std over $\geq 3$
    seeds.
  \item \textbf{Write model cards}: document intended use,
    limitations, and biases.
  \item \textbf{Publish code and data}: use GitHub and Zenodo or
    Hugging Face Datasets.
  \item \textbf{Log compute resources}: GPU type, training time,
    energy consumption.
  \item \textbf{Use deterministic algorithms}: accept the performance
    cost for critical experiments.
\end{enumerate}
\end{bestpractice}

\begin{antipattern}[title=\textbf{Reproducibility anti-patterns}]
\begin{itemize}
  \item Cherry-picking the best run out of many.
  \item ``Code available upon request''---almost never provided.
  \item Not reporting hyperparameter search ranges.
  \item Using \texttt{latest} Docker tags or unpinned dependencies.
  \item Ignoring non-determinism because ``results are approximately
    the same''.
  \item Not documenting data preprocessing steps.
\end{itemize}
\end{antipattern}

\section{Mini-project: Reproducible Experiment}
\label{sec:mini-project-ch11}

\begin{miniproject}[title=\textbf{Mini-project 11: Reproducibility}]
\begin{enumerate}
  \item Implement a complete \texttt{set\_seed()} function covering
    Python, NumPy, and PyTorch.
  \item Run the same experiment 5 times with different seeds and
    report mean $\pm$ std with 95\% confidence intervals.
  \item Write a model card for the trained model following the
    template above.
  \item Write a datasheet for the dataset used.
  \item Create a \file{README.md} with step-by-step reproduction
    instructions.
  \item Verify that a classmate can reproduce your results within
    $\pm 1\%$ by following your instructions alone.
\end{enumerate}
\end{miniproject}

\section{Exercises}
\label{sec:exercises-ch11}

\begin{exercise}[$\bigstar$ --- Seed management]
Train a small neural network on MNIST \emph{without} setting any
random seed. Run training 5 times and record the variance in test
accuracy. Then implement full seed management and repeat. Compare the
variance in both cases.
\end{exercise}

\begin{exercise}[$\bigstar\bigstar$ --- Model card creation]
Choose a pre-trained model from Hugging Face Hub. Write a detailed
model card including:
\begin{enumerate}
  \item Model details and architecture
  \item Intended use and out-of-scope uses
  \item Performance metrics disaggregated by demographic group (if
    applicable)
  \item Ethical considerations and known biases
  \item Reproduction instructions
\end{enumerate}
\end{exercise}

\begin{exercise}[$\bigstar\bigstar\bigstar$ --- Full reproducibility audit]
Take an existing ML project (your own or an open-source one) and
conduct a reproducibility audit:
\begin{enumerate}
  \item Attempt to reproduce the reported results from scratch
  \item Document every issue encountered (missing seeds, undocumented
    steps, version conflicts)
  \item Fix all reproducibility issues
  \item Add a model card, datasheet, and NeurIPS-style checklist
  \item Run multi-seed experiments and report confidence intervals
  \item Package the project with Docker for environment reproducibility
\end{enumerate}
\end{exercise}

\section{Cheatsheet}

\begin{cheatsheet}[title=\textbf{Chapter 11 Summary}]
\begin{tabular}{@{}p{4cm}p{8.5cm}@{}}
\toprule
\textbf{Concept} & \textbf{Key Point} \\
\midrule
Repeatability & Same team, same setup, same results \\
Reproducibility & Different team, same code/data, same results \\
Replicability & Different team, different code, same conclusions \\
Random seeds & Fix Python, NumPy, PyTorch, CUDA, DataLoader \\
Model card & Standardized doc: use, metrics, limitations, biases \\
Datasheet & Dataset documentation: composition, collection, uses \\
Multi-run & Report mean $\pm$ std over $\geq 3$ seeds \\
Determinism & \texttt{torch.use\_deterministic\_algorithms(True)} \\
\bottomrule
\end{tabular}
\end{cheatsheet}
